Frontier large language models exhibit a striking convergence in default responses---a phenomenon sometimes called ``persona collapse,'' where models trained on similar data produce nearly identical preferences and opinions. We investigate whether this convergence is driven by the shared Transformer architecture or by the training data itself. We train parameter-matched Transformer, GRU, and LSTM models on identical data (the \tinystories corpus) and probe their default next-word predictions, response diversity, and internal representation structure. Our experiments reveal that architecture alone produces substantially different attractor patterns: matched-capability Transformer and GRU models agree on their top-1 next-word prediction only 25\% of the time, with a mean Jensen-Shannon divergence of 0.34 between their output distributions. Transformers develop tighter representation spaces (within-concept cosine similarity of 0.94 vs.\ 0.84 for GRU) but fail to separate semantic concepts, creating a ``universal attractor basin'' that facilitates persona collapse. In contrast, recurrent architectures preserve better concept separation despite comparable perplexity. We complement these controlled experiments with a survey of frontier LLMs confirming the persona collapse phenomenon. Our results suggest that the homogeneity of default behaviors in modern LLMs is partly a consequence of shared architectural inductive biases, not solely a property of training data or alignment procedures.
